% ============================================================
%  COMSATS University Islamabad
%  MS Synopsis / Thesis – COMSATS MS Format
%  File: COMSATS-MS-Format-Thesis 3.tex
%  Compile: pdflatex "COMSATS-MS-Format-Thesis 3.tex"
%  Output:  COMSATS-MS-Format-Thesis 3.pdf
% ============================================================
\documentclass[12pt,a4paper]{article}

% ── Packages ────────────────────────────────────────────────
\usepackage[T1]{fontenc}
\usepackage[utf8]{inputenc}
\usepackage{geometry}
\geometry{top=2.5cm, bottom=2.5cm, left=3cm, right=2.5cm}
\usepackage{setspace}
\doublespacing
\usepackage{times}          % Times New Roman font
\usepackage{parskip}
\usepackage{titlesec}
\usepackage{booktabs}
\usepackage{array}
\usepackage{longtable}
\usepackage{hyperref}
\usepackage{natbib}
\usepackage{enumitem}
\usepackage{caption}
\usepackage{multirow}
\usepackage{graphicx}

% ── Section formatting ───────────────────────────────────────
\titleformat{\section}{\normalfont\large\bfseries}{\thesection.}{0.5em}{}
\titleformat{\subsection}{\normalfont\normalsize\bfseries}{\thesubsection}{0.5em}{}

% ── Hyperlink colours ────────────────────────────────────────
\hypersetup{colorlinks=true, linkcolor=black, citecolor=black, urlcolor=blue}

% ============================================================
\begin{document}

% ── Cover / Header ───────────────────────────────────────────
\begin{center}
    {\large\bfseries COMSATS University Islamabad}\\[4pt]
    {\large Synopsis for MS $\square$ \quad Ph.D.$\square$}
\end{center}

\bigskip
\noindent
\begin{tabular}{@{}p{5.5cm} p{8cm}@{}}
\textbf{Name:} Andleeb Zahra &
\textbf{Registration No.:} FA24-RBM-003 \\[4pt]
\textbf{Program:} MS Biochemistry and Molecular Biology &
\textbf{Area of Specialisation:} Biochemistry and Molecular Biology \\[4pt]
\textbf{Department:} Biosciences &
\textbf{Campus:} Islamabad \\[4pt]
\textbf{Date of Admission:} Fall 2024 &
\textbf{Date of Synopsis Submission:} October 2025 \\
\end{tabular}

\bigskip
\noindent\textbf{Proposed Title of the Thesis:}\\
Association of VDR Gene Variability, Fluoride Exposure and Fluorosis Among Children
Population of Thal Desert Areas.

\bigskip
\noindent\textbf{Supervisory Committee}

\begin{center}
\begin{tabular}{|p{8cm}|p{5cm}|}
\hline
\textbf{Name and Designation} & \textbf{Role} \\
\hline
Dr.\ Ismat Nawaz (Associate Professor) & Supervisor \\
\hline
Dr.\ Syed Ali Musstjab Akber Shah Eqani & Co-supervisor / Member \\
\hline
 & Member \\
\hline
 & Member \\
\hline
\end{tabular}
\end{center}

\vspace{1cm}
\noindent Student's Signature: \underline{\hspace{6cm}}

\bigskip\hrule\bigskip

% ── Summary ─────────────────────────────────────────────────
\noindent\textbf{Summary of the Research}

In many parts of Pakistan, especially among children, the problem of fluoride contamination in
groundwater is an unresolved issue of concern to the entire population. The constant consumption
of fluoride in excess of the standards will result in skeletal fluorosis, a disease that causes the
discolouring of enamel, malformed bones, and decreased mineralisation. Although environmental
exposure levels are similar across individuals, the degree of fluorosis manifestation varies, suggesting
that genetic diversity can affect fluoride metabolism and toxicity. The Vitamin D Receptor (VDR)
gene is one of the genetic determinants currently being studied; this gene is important in calcium
absorption and bone homeostasis. The biological response to fluoride exposure might be dependent
upon polymorphisms in this gene, especially the single-nucleotide polymorphism (SNP) rs731236
(TaqI), which alters receptor activity and downstream signalling. The current research is planned as
a cross-sectional molecular study of school-going children living in fluoride-exposure regions of
Punjab, including Thal Desert. Analyses of water and biological samples will be conducted to
determine the degree of fluoride exposure, while genotyping of the VDR TaqI SNP will be performed
using PCR-based molecular techniques. Statistical testing of genotype distribution and fluorosis
severity will be performed to establish potential gene–environment interactions. The results should
increase the body of knowledge about fluoride susceptibility at the molecular level, establish genetic
predispositions, and help design specific prevention interventions in susceptible groups.

% ── Sections ────────────────────────────────────────────────
\section{Introduction}

Fluoride is a naturally occurring element in the earth's crust and enters groundwater systems at
highly variable concentrations, reflecting regional geology and hydrochemical conditions. Within
optimum doses, fluoride improves dental health by enhancing enamel mineralisation and preventing
dental cavities \citep{WHO2017}. Excessive consumption, however, causes severe health effects
collectively known as fluorosis. Dental fluorosis results from excess fluoride during tooth
development, causing enamel defects and discolouration; the skeletal form accumulates fluoride in
bone tissue, causing pain, stiffness, and skeletal deformities.

Pakistan records many areas with groundwater fluoride above the WHO reference of 1.5\,mg/L.
Machine-learning models applied to more than 5\,000 fluoride readings estimated that approximately
13 million Pakistani residents live in regions with groundwater fluoride exceeding 1.5\,mg/L
\citep{Khan2023}. Children are particularly sensitive because developing bone and dental structures
incorporate fluoride more readily than adult tissue.

Despite similar environmental exposures, individual differences in fluorosis severity are highly
variable, indicating that genetic predisposition may be a major modifying factor \citep{Wang2022}.
The VDR gene is a strong candidate: its protein is a nuclear receptor for 1,25-dihydroxyvitamin~D
that controls calcium absorption, bone mineralisation, and skeletal homeostasis \citep{Holick2006}.
VDR polymorphisms can alter receptor expression and activity, thereby changing calcium metabolism
and the individual response to fluoride exposure. The TaqI SNP (rs731236), situated at the 3\rlap{$'$}
end of the VDR gene, is among the most investigated. Its genotypes have been associated with
changes in bone mineral density and mineral metabolism across diverse populations
\citep{Gorczynska2024}. Despite these advances, the literature examining VDR TaqI (rs731236)
and skeletal fluorosis in Pakistani children remains sparse, a gap this study aims to fill.

\section{Literature Review}

Endemic fluorosis is caused by chronic intake of high-fluoride water, brick-tea, or salt, resulting in
dose-dependent skeletal and dental phenotypes. Host genetics confer additional control over who
actually develops fluorosis, particularly among children whose enamel and bone are actively
mineralising \citep{Pramanik2017, Niazi2023}.

Early gene–environment research in high-fluoride areas demonstrated that risk is altered by
detoxification and bone-matrix genes such as \textit{GSTP1} and osteocalcin, inspiring
candidate-gene analyses in fluoride-exposed populations \citep{Wu2015}. In brick-tea regions of
China, GSTP1 polymorphisms were among the first to be linked to skeletal fluorosis, implicating
oxidative-stress pathways. Subsequent paediatric work in high-fluoride areas investigated VDR
FokI patterns in children, relating VDR genotypes to fluoride exposure conditions---an early
indication that VDR biology could determine fluoride sensitivity in development \citep{Yang2011}.

A larger cross-sectional case-control study later identified VDR FokI CT/TT genotypes as
associated with reduced risk of skeletal fluorosis in brick-tea regions \citep{Yang2016}.
Enamel-matrix genes were also found to be relevant: COL1A2 rs412777 was associated with
dental fluorosis severity, placing matrix/mineralisation genetics as central to fluoride responses
\citep{JarquinYneza2018}. A systematic review of SNPs and dental fluorosis confirmed multiple
implicated loci (matrix, signalling, antioxidant), supporting the view that genetic background
contributes substantially to inter-individual variability \citep{GonzalezCasamada2022}.

The TaqI rs731236 variant has demonstrated case-control differences in caries-focused cohorts
\citep{Antunes2016, Lei2021}, indicating that VDR signalling plays a functional role in enamel
mineral homeostasis---the same axis activated during fluoride responses in amelogenesis. A
meta-analysis of VDR polymorphisms (FokI, TaqI, ApaI, BsmI) confirmed they mediate oral
risk in children \citep{Peponis2023}, suggesting that rs731236-mediated modulation of fluoride
sensitivity is biologically plausible, although direct paediatric fluorosis data remain limited.

More recent studies continue to implicate alternative genes: COL1A2 variants were linked to
dental fluorosis in Tunisian and Latin-American cohorts \citep{Kallala2024, JarquinYneza2018},
while AQP5, MMP2, and COMT polymorphisms were investigated in Brazilian endemic fluorosis
populations \citep{deLavor2025}. These findings support a polygenic model of fluoride
susceptibility. Thai school children also revealed salivary proteome changes in severe fluorosis
\citep{Gavila2024}, and a 2024 study detected genetic instability signals in fluorosis patients,
confirming systemic cellular stress from chronic fluoride exposure \citep{Garcia2024}.

A BMC Pediatrics meta-analysis (2024) confirmed the association of rs731236 with caries risk in
children and suggested that the 3\rlap{$'$}-end position of this SNP may interfere with mRNA
processing or VDR structure, affecting functional signalling in mineralising tissues \citep{Qin2024}.
AI-assisted models now incorporate both genotype and environmental data to predict enamel outcomes
\citep{Wei2025}, a framework that could be adapted to fluorosis risk prediction with inclusion of VDR.

In summary, while VDR FokI remains the most directly studied locus in fluorosis, the biological
plausibility of rs731236 (TaqI) involvement is well-supported by parallel caries and bone-density
literature. Child-specific, properly powered studies combining TaqMan/NGS genotyping, urinary
fluoride measurement, 25(OH)D assessment, and standardised fluorosis indices (TFI/Dean's) are
needed to determine whether rs731236 stratifies fluorosis risk/severity in endemic settings.

\section{Problem Statement}

High fluoride levels in drinking water remain a significant public health concern in Pakistan, where
groundwater fluoride frequently exceeds the WHO guideline of 1.5\,mg/L. In heavily contaminated
regions such as Punjab, dental and skeletal fluorosis have become endemic, manifesting as enamel
mottling, defective bone mineralisation, chronic pain, and skeletal deformities. Although fluoride
exposure is the primary environmental driver, considerable inter-individual variation in fluorosis
incidence and severity exists among people exposed to comparable fluoride levels, implying a genetic
component in the response to fluoride toxicity.

Emerging evidence suggests that polymorphisms in the VDR gene---a key regulator of calcium
absorption, bone remodelling, and mineral homeostasis---may modulate fluorosis susceptibility.
Among these, the VDR TaqI (rs731236) SNP has been hypothesised to influence receptor activity
and downstream calcium metabolism, thereby affecting the skeletal response to chronic fluoride
exposure. Despite global interest in fluoride–gene interactions, the role of VDR rs731236 in
skeletal fluorosis remains understudied, particularly in Pakistan.

This study therefore aims to examine the relationship between VDR TaqI (rs731236) polymorphism
and fluorosis severity in children residing in high-fluoride regions of Punjab, combining
environmental exposure data with genetic profiling to identify genetically vulnerable subgroups.
Findings will inform targeted public health measures and advance understanding of gene–environment
interactions underlying fluoride toxicity in vulnerable Pakistani populations.

\section{Objectives}

\begin{enumerate}[label=\textbullet]
\item To evaluate fluoride levels in water and to determine the incidence of skeletal fluorosis in
      children in the region of Punjab, Pakistan.
\item To analyse samples at the genetic level for potential single-nucleotide polymorphisms (SNPs)
      using tetra-primer ARMS-PCR.
\item To identify the relationship among fluoride exposure, skeletal fluorosis, and single-nucleotide
      polymorphisms (SNPs).
\end{enumerate}

\section{Research Methodology}

\subsection{Ethical Consideration}

Before initiating the research, formal approval will be sought from the Ethical Review Board (ERB)
of the Department of Biosciences, COMSATS University Islamabad. The study will be conducted in
the Thal Desert area of Punjab, which has documented high groundwater fluoride levels.

\subsection{Population and Sampling Strategy}

One hundred children aged 12--16 years (Grades 7--10) will be recruited from schools in the study
area using stratified random sampling. Each participant will provide a 5\,mL venous blood sample:
2\,mL in EDTA vacutainers for DNA extraction and 3\,mL in gel vacutainers for serum separation.
Serum samples will be centrifuged and stored at $-20\,^\circ$C pending biochemical analysis.

\subsection{Fluoride Quantification}

Fluoride concentrations in drinking water and serum samples will be measured by the
ion-selective electrode (ISE) method to obtain true individual exposure levels.

\subsection{Genomic DNA Isolation}

Genomic DNA will be extracted using the Phenol--Chloroform--Isoamyl Alcohol (PCI) method.
Extracted DNA will be assessed by 2\% agarose gel electrophoresis and stored at $-20\,^\circ$C
until further molecular analysis. Specific primers for the VDR TaqI (rs731236) polymorphism
will be designed using Primer3 software and validated for specificity, efficiency, and amplification
consistency.

\subsection{PCR and Gel Analysis / Genotyping}

VDR TaqI polymorphism will be genotyped by Tetra-Primer Amplification Refractory Mutation
System PCR (Tetra-ARMS PCR). Amplified products will be resolved by agarose gel
electrophoresis; homozygous and heterozygous genotypes will be distinguished by fragment size.

\subsection{Statistical Analysis}

Data will be entered into SPSS (version 22) and GraphPad Prism. Hardy--Weinberg equilibrium
will be assessed for VDR rs731236 genotype frequencies. Comparative genotypic and allelic
frequencies between fluorosis and control groups will be used to measure associations with disease
susceptibility, with statistical significance set at $p = 0.05$. Logistic regression will be applied to
identify independent relationships among VDR polymorphism, fluoride exposure, and fluorosis
severity.

% ── Tentative Schedule ───────────────────────────────────────
\section*{Tentative Schedule}

\begin{center}
\begin{tabular}{|p{4.5cm}|c|c|c|c|c|c|c|c|}
\hline
\textbf{Activity} &
\rotatebox{90}{Sept--Oct} &
\rotatebox{90}{November} &
\rotatebox{90}{December} &
\rotatebox{90}{January} &
\rotatebox{90}{February} &
\rotatebox{90}{March} &
\rotatebox{90}{April} &
\rotatebox{90}{May--June} \\
\hline
Literature review \& Ethical Approval &\checkmark& & & & & & & \\
\hline
Sampling & &\checkmark&\checkmark& & & & & \\
\hline
DNA Extraction \& Genotyping & & &\checkmark&\checkmark& & & & \\
\hline
Fluoride Quantification & & &\checkmark&\checkmark& & & & \\
\hline
Statistical Analysis & & & & &\checkmark&\checkmark& & \\
\hline
Thesis Writing & & & & & &\checkmark&\checkmark& \\
\hline
Submission & & & & & & & &\checkmark \\
\hline
\end{tabular}
\end{center}

% ── References ───────────────────────────────────────────────
\bibliographystyle{apalike}
\begin{thebibliography}{99}

\bibitem[Antunes et~al.(2016)]{Antunes2016}
Antunes, L.\ A.\ A., et~al.\ (2016).
VDR TaqI (rs731236) polymorphism in children with and without caries.
\textit{Caries Research}, 51(1), 7--13.

\bibitem[Chisini et~al.(2025)]{Chisini2025}
Chisini, L.\ A., et~al.\ (2025).
Pathways of VDR gene and dental caries: A systematic review and meta-analysis.
\textit{Archives of Oral Biology}.

\bibitem[de~Lavôr et~al.(2025)]{deLavor2025}
de~Lavôr, J.\ R., et~al.\ (2025).
AQP5, MMP2, and COMT polymorphisms in children from endemic fluorosis areas.
\textit{Caries Research}.

\bibitem[García et~al.(2024)]{Garcia2024}
García, A.\ L.\ H., et~al.\ (2024).
Genetic instability in dental fluorosis.
\textit{Science of the Total Environment}, 906, 67393.

\bibitem[Gavila et~al.(2024)]{Gavila2024}
Gavila, P., et~al.\ (2024).
Salivary proteomic signatures in severe dental fluorosis in children.
\textit{Scientific Reports}, 14, 69409.

\bibitem[González-Casamada et~al.(2022)]{GonzalezCasamada2022}
González-Casamada, C., et~al.\ (2022).
SNPs and dental fluorosis: A systematic review.
\textit{Dentistry Journal}, 10(11), 211.

\bibitem[Górczyńska-Kosiorz et~al.(2024)]{Gorczynska2024}
Górczyńska-Kosiorz, S., Tabor, E., Niemiec, P., Pluskiewicz, W., \& Gumprecht, J.\ (2024).
Associations between the VDR gene rs731236 (TaqI) polymorphism and bone mineral density in
postmenopausal women from the RAC-OST-POL.
\textit{Biomedicines}, 12(4), 917.

\bibitem[Holick(2006)]{Holick2006}
Holick, M.\ F.\ (2006).
Vitamin D: Its role in bone health and disease.
\textit{Annals of Internal Medicine}.

\bibitem[Jarquín-Yñezá et~al.(2018)]{JarquinYneza2018}
Jarquín-Yñezá, L., et~al.\ (2018).
Dental fluorosis and COL1A2 rs412777.
\textit{Archives of Oral Biology}.

\bibitem[Kallala et~al.(2024)]{Kallala2024}
Kallala, R., et~al.\ (2024).
Association between dental fluorosis and COL1A2 rs412777.
\textit{BMC Oral Health}, 24, 4086.

\bibitem[Khan et~al.(2023)]{Khan2023}
Khan, S., et~al.\ (2023).
Machine-learning modelling of groundwater fluoride in Pakistan.
\textit{Environmental Science \& Technology}.

\bibitem[Lei et~al.(2021)]{Lei2021}
Lei, W., et~al.\ (2021).
TaqI (rs731236 T$>$C) and caries.
\textit{Caries Research}, 55(3).

\bibitem[Niazi \& Al-Khanati(2023)]{Niazi2023}
Niazi, F.\ C., \& Al-Khanati, N.\ (2023).
Dental fluorosis.
\textit{StatPearls} (Updated 2023). NCBI.

\bibitem[Peponis et~al.(2023)]{Peponis2023}
Peponis, M., et~al.\ (2023).
Vitamin D/VDR polymorphisms and oral outcomes: A meta-analysis.
\textit{Applied Sciences}, 13(10), 6014.

\bibitem[Pramanik \& Saha(2017)]{Pramanik2017}
Pramanik, S., \& Saha, D.\ (2017).
The genetic influence in fluorosis.
\textit{Mutation Research/Reviews}.

\bibitem[Qin et~al.(2024)]{Qin2024}
Qin, X., et~al.\ (2024).
VDR polymorphisms and caries risk in children: A meta-analysis.
\textit{BMC Pediatrics}, 24, 5127.

\bibitem[Wang et~al.(2022)]{Wang2022}
Wang, X., et~al.\ (2022).
Dental fluorosis presents with differing severity in individuals exposed to similar levels of
fluoride intake, suggesting the influence of genetic factors.
\textit{Frontiers in Genetics}.

\bibitem[Wei et~al.(2025)]{Wei2025}
Wei, J., et~al.\ (2025).
AI-assisted models for enamel outcome prediction under fluoride and genotype conditions.
\textit{Environmental Health Perspectives}.

\bibitem[Wen et~al.(2022)]{Wen2022}
Wen, C., et~al.\ (2022).
Brick-tea consumption and fluorosis: A review.
\textit{Frontiers in Nutrition}, 9, 1030344.

\bibitem[WHO(2017)]{WHO2017}
World Health Organization.\ (2017).
\textit{Guidelines for Drinking-water Quality} (4th ed., incorporating 1st addendum).
WHO Press.

\bibitem[Wu et~al.(2015)]{Wu2015}
Wu, J., Wang, W., Liu, Y., Sun, J., Ye, Y., Li, B., \ldots{} \& Gao, Y.\ (2015).
Modifying role of GSTP1 polymorphism on the association between tea fluoride exposure and
brick-tea type fluorosis.
\textit{PLOS ONE}, 10(6), e0128280.

\bibitem[Yang et~al.(2011)]{Yang2011}
Yang, Y., et~al.\ (2011).
Relationship between VDR FokI genotypes in children from high- and low-fluoride areas.
\textit{Chinese Journal of Public Health}.

\bibitem[Yang et~al.(2016)]{Yang2016}
Yang, D., et~al.\ (2016).
Association between VDR FokI polymorphism and skeletal fluorosis of brick-tea type:
A case-control study.
\textit{BMJ Open}, 6, e011980.

\bibitem[Zhang et~al.(2023)]{Zhang2023}
Zhang, K., et~al.\ (2023).
Epidemiology and pathogenesis of fluorosis: Genetics matter.
\textit{Frontiers in Cell and Developmental Biology}, 11, 1168215.

\end{thebibliography}

\end{document}
